\section{A simultaneous weighted prime race}\label{sec:weightedrace}

Chebyshev's original exponential smoothing asks whether
\[
   \sum_{p>2}(-1)^{(p-1)/2} e^{-p/x}\;\longrightarrow\;-\infty .
\]
Hardy--Littlewood and Landau identified this assertion with the zero
location for the nonprincipal character modulo four. The repository's theorem supplies
that input. For the precise equivalence and its history see M.~Suzuki,
\emph{On variants of Chebyshev's conjecture}, \S1.2, equations (16)--(17),
\href{https://doi.org/10.1007/s11139-025-01238-9}{doi:10.1007/s11139-025-01238-9}.
With a different smoothing we can go beyond this single race and
measure the lead in every nonsquare class. The following deduction keeps track of
possible zeros at the center.

\paragraph{Theorem.}
Fix a modulus $q$, a quotient $G$ of $(\mathbb Z/q\mathbb Z)^\times$,
and write $h=|G|$, $G^2=\{a^2:a\in G\}$, $G[2]=\{a:a^2=1\}$.
For $g\in G$ put
\[
 S_g(x)=\sum_{\substack{p\le x,\;p\nmid q\\ {[p]}=g}}
            \frac{\log p}{\sqrt p}\log\frac{x}{p}.
\]
For every nonprincipal $\chi\in\widehat G$ let
$m_\chi=\operatorname{ord}_{s=1/2} L(s,\chi)$, using the character
pulled back to modulus $q$. Then, simultaneously for all $g$,
\[
 S_g(x)-S_1(x)=B_g(\log x)^2+O_q(\log x),
\]
where
\[
 B_g=\frac{|G[2]|}{4h}\,\mathbf 1_{g\notin G^2}
    +\frac1{2h}\sum_{\chi\ne1}
       (1-\operatorname{Re}\chi(g))m_\chi .
\]
In particular, every nonsquare coset eventually leads the identity coset,
by at least $\bigl(|G[2]|/(4h)+o(1)\bigr)(\log x)^2$.
If every element of $G$ has order dividing two, all $g\ne1$ lead by
at least $(1/4+o(1))(\log x)^2$. For a multiquadratic extension of
$\mathbb Q$, the identity class describes primes that split completely,
and the other cosets specify patterns of quadratic splitting at unramified primes.

\paragraph{Proof.}
For a nonprincipal character let $u=\log x$. Suzuki's explicit formulas
(62) and (64), in the proof of Theorem~8 of the cited paper, give
\[
 A_\chi(x)=\sum_{n\le x}\frac{\Lambda(n)\chi(n)}{\sqrt n}\log(x/n)
      =-\frac{m_\chi}{2}u^2+O_q(u).
\]
Indeed the central zero gives the quadratic term. The remaining nontrivial
zeros $1/2+i\gamma$ contribute a bounded series, since
$|x^{i\gamma}|=1$ by Dirichlet GRH and
$\sum_{\gamma\ne0}|\gamma|^{-2}<\infty$ for each fixed character,
with multiplicities. All other terms are $O_q(u)$.
Passing from a primitive character to a character it induces costs
$O_q(u)$ in $A_\chi$, by summing the geometric tails at the finitely many
removed primes. The extra Euler factors do not vanish at $s=1/2$.

The terms $n=p^k$, $k\ge3$, contribute $O(u)$ because
$\sum_p (\log p)p^{-3/2}/(1-p^{-1/2})$ converges. The terms $n=p^2$ equal
\[
 2\int_1^{\sqrt x}
        \left(\sum_{p\le t}\frac{\chi^2(p)\log p}{p}\right)\frac{dt}{t}
 =\mathbf 1_{\chi^2=1}\frac{u^2}{4}+O_q(u).
\]
For principal $\chi^2$ the inner sum is $\log t+O_q(1)$ by Mertens'
estimate. For nonprincipal $\chi^2$ it is $O_q(1)$ by partial summation
of the prime number theorem in progressions at fixed modulus, in its
classical form with error $O_q(t\exp(-c_q\sqrt{\log t}))$.
Subtracting the prime powers thus gives a prime-only sum with quadratic
coefficient $-m_\chi/2-\mathbf 1_{\chi^2=1}/4$.

Character orthogonality expresses $S_g-S_1$ as the sum of those
prime-only sums with coefficients $(\overline{\chi(g)}-1)/h$. The principal
character cancels. Conjugate characters have equal central orders, so the
final coefficient may be written using real parts. Characters satisfying
$\chi^2=1$ are the dual of $G/G^2$. Their sum at $g$ equals
$|G/G^2|$ for $g\in G^2$ and zero otherwise. Since
$|G/G^2|=|G[2]|$, this gives the stated coefficient. Each central-zero
term is nonnegative. The group is finite, so the errors can be taken
simultaneously for all classes.\hfill$\square$

\paragraph{Example modulo eight.}
For each $g=3,5,7$,
\[
  S_g(x)-S_1(x)=\tfrac14(\log x)^2+O(\log x).
\]
Here all three central orders are zero, as one can see without a
numerical calculation. On the odd classes $1,3,5,7$ the three respective
sign patterns are $+,-,+,-$, $+,+,-,-$, and $+,-,-,+$.
The Dirichlet series at $s=1/2$ converges by periodicity and zero mean.
Each consecutive block of four odd summands has positive sum for all
three patterns. For the last pattern this uses
$f(a)-f(a+2)>f(a+4)-f(a+6)$ with $a=8k+1$ and
$f(t)=t^{-1/2}$, which follows from strict convexity. Thus every central
value is positive. The example gives three strict eventual inequalities
with a common leading coefficient. The theorem concerns the specified
weighted sums; an unweighted two-way prime race can change sign.

\paragraph{Provenance.}
Suzuki's Theorems~3 and~4 already give modulo-four conclusions and his
Theorem~9 gives the individual real-character prime sum for a nonzero
central value. The quotient formula above combines his explicit formula
with the square map on $G$; it is presented as a deduction without a claim
of priority.
